For an open $U\subset X$ and an open $V\supset f(U)$, restriction gives $\mathscr F(f^{-1}V)\to\mathscr F(U)$. These maps induce a morphism from the presheaf defining $f^{-1}f_*\mathscr F$ to $\mathscr F$, hence a natural morphism $\varepsilon:f^{-1}f_*\mathscr F\to\mathscr F$.

For an open $V\subset Y$, a section of $\mathscr G(V)$ represents a section of the inverse-image presheaf on $f^{-1}V$. Composing with sheafification gives $\mathscr G(V)\to(f^{-1}\mathscr G)(f^{-1}V)$, hence a natural morphism $\eta:\mathscr G\to f_*f^{-1}\mathscr G$.

The proposed maps between Hom sets are $\beta\mapsto f_*\beta\circ\eta$ and $\alpha\mapsto\varepsilon\circ f^{-1}\alpha$. Explicitly, the second sends a local representative $s\in\mathscr G(V)$ of an inverse-image section on $U\subset f^{-1}V$ to $\alpha(V)(s)|_U$. Starting with $\alpha$ therefore recovers its value on every section of every $V$. Starting with $\beta$ recovers it on all these local representatives, which generate the inverse-image sheaf locally, and hence recovers $\beta$. The constructions commute with morphisms in either sheaf, so the bijection is natural.
